% ===================================================================
\chapter{Introduction}
\label{ch:intro}

\section{The problem}
\label{sec:problem}

A modern high-resolution automotive radar generates data faster than the
vehicle it is mounted on can conveniently move it. The cascaded radar used in
this thesis produces about \SI{3}{\mebi\byte} of raw digitised samples every
frame, and it is one of several such sensors on a car. Getting that data to the
central processor means either expensive shielded high-speed cabling or
compressing it at the sensor first. Compression is the cheaper answer, but it
has to keep up with the sensor, fit in a small amount of silicon, and --- because
this data feeds a system that decides whether to brake --- lose nothing. This
thesis compares the available compression methods on real radar data, selects
the ones that meet those conditions, and builds and measures them on an FPGA.

\section{Motivation}
\label{sec:motivation}

Radar earns its place in a vehicle's sensor suite because it works when the
other sensors do not. Cameras fail in fog, glare and darkness; lidar is
degraded by rain and airborne dust; radar sees through all of them
\cite{waldschmidt2021,yao2023}. That robustness is why radar has been in
production vehicles for two decades, initially for simple distance warning and
now for adaptive cruise control, automatic emergency braking, blind-spot
detection and lane-change assistance \cite{skolnik2008}.

What has changed recently is resolution. Supporting SAE Level 2+ autonomy needs
more than ``something is \SI{40}{\metre} ahead''. It needs a map: how many
objects, where, how fast, how wide, and what they probably are. Angular
resolution is what separates ``an object'' from ``a car and a pedestrian
standing beside it'', and angular resolution comes from having many antennas.
A radar with $N_t$ transmitters and $N_r$ receivers can synthesise an effective
array of $N_t \times N_r$ elements, so resolution improves roughly with the
product. The data volume, unfortunately, scales the same way.

At the same time, vehicle electrical architectures have been consolidating.
Where a car once had a hundred small electronic control units each doing one
job, the industry has moved to domain controllers, then to \emph{zonal}
architectures where sensors report to a nearby gateway, and toward fully
centralised designs where one high-performance system-on-chip does everything
\cite{bandur2021,liu2022,barabas2024}. Centralising the computation has obvious
advantages, but it means the raw sensor data has to travel.

Figure~\ref{fig:arch} shows the consequence. A Serialiser/Deserialiser (SerDes)
link \cite{ti_serdes,gmsl} carries the aggregated output of several radar
modules at \SI{2.5}{Gbit\per\second} or more. Those data rates are achievable,
but they require shielded cable, they cost money and weight, and they put a
matching ingest burden on the SoC's memory system \cite{matheus2021}. If
instead the data is compressed at the radar module by a factor of \num{2.5} or
better, the same information fits inside Gigabit Ethernet
\cite{ieee8023bp} --- a cheap, standard, unshielded interconnect. That single
substitution is the economic argument for this entire body of work.

\begin{figure}[t]
\centering
\begin{tikzpicture}[font=\small,
  box/.style={draw=slate,fill=white,rounded corners=2pt,minimum height=7mm,inner xsep=4pt},
  rad/.style={box,fill=lightbg,minimum width=13mm},
  lbl/.style={font=\scriptsize,text=slate}]

% ---- top: without compression
\node[lbl,anchor=west] at (-0.2,1.9) {\textbf{(a) Raw transport}};
\node[rad] (r1) at (0.9,1.1) {radar};
\node[rad] (r2) at (0.9,0.2) {radar};
\node[box,minimum width=14mm] (ser) at (4.0,0.65) {serialiser};
\node[box,minimum width=14mm] (des) at (8.6,0.65) {deserialiser};
\node[box,minimum width=12mm,fill=lightbg] (soc) at (11.9,0.65) {SoC};
\draw[-{Latex[length=1.6mm]},slate] (r1) -- (ser);
\draw[-{Latex[length=1.6mm]},slate] (r2) -- (ser);
\draw[-{Latex[length=1.6mm]},warm,line width=1.4pt] (ser) -- node[above,lbl,text=warm,yshift=1pt]{\SI{2.5}{Gbit/s} shielded serial} (des);
\draw[-{Latex[length=1.6mm]},slate] (des) -- (soc);

% ---- bottom: with compression
\node[lbl,anchor=west] at (-0.2,-0.75) {\textbf{(b) Compressed transport}};
\node[rad] (q1) at (0.9,-1.55) {radar};
\node[rad] (q2) at (0.9,-2.45) {radar};
\node[box,minimum width=16mm,fill=lightbg] (cmp) at (4.0,-2.0) {compressor};
\node[box,minimum width=18mm,fill=lightbg] (dec) at (8.8,-2.0) {decompressor};
\node[box,minimum width=12mm,fill=lightbg] (soc2) at (11.9,-2.0) {SoC};
\draw[-{Latex[length=1.6mm]},slate] (q1) -- (cmp);
\draw[-{Latex[length=1.6mm]},slate] (q2) -- (cmp);
\draw[-{Latex[length=1.6mm]},green2,line width=1.4pt] (cmp) -- node[above,lbl,text=green2,yshift=1pt]{$<$\,\SI{1}{Gbit/s} Gigabit Ethernet} (dec);
\draw[-{Latex[length=1.6mm]},slate] (dec) -- (soc2);
\end{tikzpicture}
\caption[Raw versus compressed radar transport]{The economic argument for
sensor-edge compression. Compressing by a factor of \num{2.5} or better at the
radar module replaces an expensive shielded high-speed serial link with
standard Gigabit Ethernet, without changing what the central processor
receives.}
\label{fig:arch}
\end{figure}

\section{Why the compressor must be lossless}
\label{sec:lossless}

It is tempting to relax the requirement. Lossy compression reaches much higher
ratios, and radar data is noisy anyway, so surely a little extra noise is
harmless?

It is not, and the reason is specific rather than philosophical. The
compressor sits \emph{upstream} of detection. Whatever it discards is discarded
before the detector ever sees it, and the detector's job is to find weak
returns close to the noise floor --- a pedestrian at range, a motorcycle between
two trucks. Lossy compression raises the noise floor precisely where the margin
is thinnest. Chapter~\ref{ch:results} quantifies this on real data: Block
Adaptive Quantization reaches the same compression ratio as the lossless coders
studied here but misses nearly three times as many real detections and reports
nearly nine times as many false ones.

Losslessness also makes verification unambiguous. The pass criterion is that
every reconstructed sample equals its original exactly. There is no tolerance
to argue about and no distortion metric to trade off: a single differing bit
anywhere in the \num{157} million input bits of the test set is a failure, and
the test reports it as one. Every design in this thesis is checked against that
criterion, from MATLAB through C-simulation and RTL co-simulation to the
encoders and decoders running on the FPGA board.

\section{Aim and objectives}
\label{sec:objectives}

The starting point is the master's thesis of Lukas Kiem \cite{kiem2025}, which
evaluated a range of radar compression schemes on recordings from a
single-transmitter Infineon radar, selected Doppler Redundancy Hybrid Encoding
(DRHE) as the best of them, and implemented its compressor on a Zynq
UltraScale+ device (compression ratio \num{3.17} in hardware). Chapter~\ref{ch:kiem}
describes that work. Four questions that his work leaves open define the
objectives of this thesis.

\begin{keybox}
\textbf{Aim.} Compare compression methods side by side on real, public,
multi-transmitter radar data; select the best \emph{lossless} methods; and
implement both their encoders and their decoders on an FPGA, measured on a
real board.
\end{keybox}

\begin{enumerate}
\item \textbf{Validate on public MIMO data.} Kiem's results come from private
recordings of a single-transmitter radar. This thesis uses 50 consecutive
frames of the public ColoRadar dataset \cite{kramer2022}, recorded with a Texas
Instruments twelve-transmitter, sixteen-receiver cascade, for every result
from simulation to hardware.

\item \textbf{Examine the quantisation step.} Kiem identified the conversion to
16-bit fixed point (FX16) as the source of his detection errors. This thesis
measures what that step costs in detection, and whether floating-point storage
is a usable alternative.

\item \textbf{Broaden the comparison.} Eight methods are compared on the same
data with the same metrics: two reference formats, three lossy quantisers, DRHE
on integer and on floating-point data, and a linear predictive coder (LPC)
designed for this thesis as a controlled comparison with DRHE.

\item \textbf{Implement and measure the complete pipeline.} Kiem implemented
compression only and measured it on his own system. This thesis implements
compression \emph{and} decompression for both selected algorithms, verifies them
bit-exactly from MATLAB to the board, and measures throughput, latency and
resources on an AMD KCU116 board.
\end{enumerate}

\section{Contributions}
\label{sec:contributions}

\begin{enumerate}
\item \textbf{A comparison of eight compression methods on real cascade radar
data} (Section~\ref{sec:stage1}). Only the two predictive lossless coders, DRHE
(\num{3.75}) and LPC (\num{3.70}), compress by more than three without losing
detections; the lossy quantisers reach up to \num{3.76} but miss
\SIrange{10}{29}{\percent} of targets, and floating-point storage compresses by
only \num{1.04}.

\item \textbf{The correct prediction axis for time-division MIMO radar}
(Section~\ref{sec:predaxis}). On a TDM cascade consecutive ramps come from
different transmitters, so slow-time prediction must use the previous ramp of
the same transmitter, $\nTx$ ramps back. Predicting from the adjacent ramp, as
single-transmitter descriptions of these algorithms do, compresses worse than not
predicting at all.

\item \textbf{A breakdown of where the compression comes from}
(Sections~\ref{sec:decomp} and~\ref{sec:ratiomeaning}). Huffman coding alone
achieves \num{3.48} on this data --- mostly from unused fixed-point headroom and
the sparsity of a range-FFT frame --- prediction adds \SIrange{6}{8}{\percent},
and the Huffman dictionary decides the rest.

\item \textbf{Huffman dictionaries matched to the prediction errors}
(Section~\ref{sec:dictresult}). Re-derived from five frames, they raise DRHE to
\num{4.22} and LPC to \num{4.40}, hold on the 45 frames they were not trained
on, and reverse the ranking of the two algorithms.

\item \textbf{Complete, verified hardware for both algorithms}
(Chapters~\ref{ch:hardware} and~\ref{ch:results}). Encoders and decoders were
implemented in Vitis HLS and Vivado and run in a board test system on an AMD
KCU116. Compressed output is byte-identical to the simulation on all 50 frames
with both dictionaries, the hardware decoders reconstruct every frame exactly,
and throughput, latency and resources are measured for both halves.

\item \textbf{A comparison with the baseline} (Chapters~\ref{ch:hwinterpret} and~\ref{ch:vskiem}),
including a reconstruction of Kiem's simulation from his text, which reaches
his simulated DRHE ratio at one plausible noise level and explains how his
results relate to the ones here.
\end{enumerate}

\section{Thesis outline}
\label{sec:structure}

Chapter~\ref{ch:background} gives the background on FMCW and MIMO radar,
lossless compression and FPGAs, and Chapter~\ref{ch:kiem} summarises Kiem's
thesis. Chapter~\ref{ch:methods} describes the eight methods and the two
predictive coders in detail, and Chapter~\ref{ch:evaluation} how they are
evaluated and verified. Chapter~\ref{ch:hardware} describes the encoders,
decoders and board test system, and Chapter~\ref{ch:results} reports the
results. Chapters~\ref{ch:hwinterpret} and~\ref{ch:vskiem} compare the hardware
and the compression results with Kiem's. Chapter~\ref{ch:discussion} discusses
what the results mean for a radar sensor, and Chapter~\ref{ch:conclusion}
concludes.
